\documentclass[10pt,a4paper]{amsart}
\usepackage[margin=19mm]{geometry}
\usepackage{amsmath,amssymb,mathtools}
\usepackage[scr=rsfs,cal=euler]{mathalfa}
\usepackage{xfrac}
\usepackage[unicode,hidelinks,bookmarks=false]{hyperref}
\newcommand{\ie}{i.e.\ }
\newcommand{\cf}{cf.\ }
\newcommand{\resp}{respectively\ }
\newcommand{\Subsection}{\subsection}
\providecommand{\nobreakdash}{\nobreak\hbox{-}}
\setlength{\emergencystretch}{2em}
\multlinegap0pt
\usepackage{amssymb}
\newtheorem{proposition}{Proposition}
\newcommand{\N}{\mathbb{N}}
\DeclareMathOperator{\Perm}{\mbox{Perm}}
\newcommand{\Z}{\mathbb{Z}}
\title{Bicyclic biskew braces}
\author{Hal Simpson, Paul J. Truman}
\date{Source: arXiv:2606.23004v1 (CC BY 4.0)}
\begin{document}
\maketitle

\noindent\emph{Source and licence.} Bicyclic biskew braces. Version arXiv:2606.23004v1 (CC BY 4.0), distributed by arXiv under the Creative Commons Attribution 4.0 International licence, http://creativecommons.org/licenses/by/4.0/, as recorded in the arXiv licence metadata for this exact version and shown on the abstract page https://arxiv.org/abs/2606.23004v1. Full licence notice: \texttt{licenses/arxiv-2606.23004-CC-BY-4.0.txt}.\par\smallskip

\noindent\textit{Editorial scope.} Counted unit: the article's proposition prop\_biskew\_general (source0.tex lines 380--382). The article's complete original proof of it is delivered with it (source0.tex lines 383--399, one \texttt{proof} environment of 17 lines). No mathematics is written, repaired or re-derived: the statement and the proof are the article's own lines, in the article's own notation, and no retained line is substituted or re-flowed.  Retained uncounted as the source-local context the proof needs: lines 97--97; lines 368--370.  Nothing was omitted from any retained span, so the package carries no omission record.  No retained line contains a citation, so the wrapper carries no bibliography and no bibliography entry.  No retained line contains a figure environment, an image-inclusion command or drawing code, so no image is delivered and the package declares no asset dependency.  Editorial formatting only, in addition: an AMS \texttt{amsart} preamble that restates the article's own declarations and macros used by the retained lines, namely the environments \texttt{proposition} on separate counters, together with the standard packages those lines need.  The retained environments print in this document as Proposition 1 in order of appearance, whereas the article prints the same items with the numbering of its own full document.  The complete native source of the article as distributed in the arXiv e-print for this exact version is bundled unchanged as \texttt{sources/203400/source0.tex}.
\section{Tools for classifying skew braces and biskew braces} \label{sec_tools}

\begin{equation} \label{eqn_circle}
i \circ_{w} j = i + j + wij
\end{equation}

\begin{proposition} \label{prop_biskew_general}
Let $ n \in \N $ and $ w $ be a divisor of $ n $ such that the binary operation $ \circ $ defined by \eqref{eqn_circle} makes $ (\Z_{n},+,\circ) $ into a bicyclic brace. Then $ (\Z_{n},+,\circ) $ is a biskew brace if and only if $ n \mid w^{2} $.  
\end{proposition}

\begin{proof}
Recall from Section \ref{sec_tools} that $ (\Z_{n},+,\circ) $ is a biskew brace if and only if the image of the left regular representation $ \lambda_{+} : \Z_{n} \rightarrow \Perm(\Z_{n}) $ normalizes the image of the left regular representation $ \lambda_{\circ} : \Z_{n} \rightarrow \Perm(\Z_{n}) $. Since $ (\Z_{n},+) $ and $ (\Z_{n},\circ) $ are both cyclic and generated by $ 1 $, this occurs if and only if $ \lambda_{+}(1) \lambda_{\circ}(1) \lambda_{+}(1)^{-1} \in \lambda_{\circ}(\Z_{n}) $. That is, if and only if 
\begin{equation} \label{eqn_biskew_condition}
\lambda_{+}(1) \lambda_{\circ}(1) \lambda_{+}(-1) = \lambda_{\circ}(j) \mbox{ for some } j.
\end{equation}
For $ i \in \Z_{n} $ we have
\begin{eqnarray*}
\lambda_{+}(1) \lambda_{\circ}(1) \lambda_{+}(-1) [i] & = & 1+( 1 \circ (i-1) ) \\
& = & 1 + 1+(i-1)+w(i-1) \\
& = & 1+i+w(i-1),
\end{eqnarray*}
whereas $ \lambda_{\circ}(j)[i] = i+j+wij $. Thus \eqref{eqn_biskew_condition} holds if and only if 
\[ 1+i+w(i-1) \equiv i+j+wij \pmod{n} \mbox{ for all } i. \]
Choosing $ i=0 $ we see that $ j=1-w $, so the condition becomes 
\[ 1+i+w(i-1) \equiv i+1-w+wi(1-w) \pmod{n} \mbox{ for all } i; \]
that is, $ w^{2} \equiv 0 \pmod{n} $. 
\end{proof}

\end{document}
